\section*{Chapitre \ref{ch:iv1:random-walks-and-martingales} --- Marches aléatoires et martingales}
\markboth{\protect\omTocMark{}{Réponses~: Marches aléatoires et martingales}}{\protect\omTocMark{}{Réponses~: Marches aléatoires et martingales}}

\begin{solution}{iq:iv1:random-walks-and-martingales:1}
On termine à 10 avec une probabilité de $0.3$, et on fait 21 pas en moyenne.

On utilise l'arrêt optionnel\index{arrêt optionnel}. La position est une martingale\index{martingale}, un processus dont l'espérance de la valeur suivante est la valeur courante, si bien que son espérance finale égale le point de départ~: $10p = 3$, donc $p = 0.3$.

Pour le temps, $S_t^2 - t$ est aussi une martingale, car un pas ajoute en moyenne 1 au carré. À la fin, $E[S_T^2] - E[T] = 3^2$, et $E[S_T^2] = 0.3 \cdot 100 = 30$. Donc $E[T] = 30 - 9 = 21$, ce qui correspond au cas général $k(n - k)$.

\iqwrong{Résoudre neuf équations linéaires, ce qui est juste mais lent. Utiliser les deux martingales évite ce travail, et elles se généralisent à des questions plus difficiles.}
\end{solution}

\begin{solution}{iq:iv1:random-walks-and-martingales:2}
On atteint 10~\$ avec une probabilité d'environ $0.884$, et il faut environ $19.2$ paris en moyenne.

Avec une pièce biaisée, $r^{S_t}$ est une martingale\index{martingale exponentielle}, où $r = q/p = 2/3$~: sur un pas, $p r + q/r = q + p = 1$. En partant de 5, et en notant $P$ la probabilité d'atteindre 10, l'arrêt optionnel donne
\[
r^5 = P r^{10} + (1 - P) \quad\text{donc}\quad P = \frac{1 - r^5}{1 - r^{10}} = \frac{1}{1 + r^5} \approx 0.8836.
\]
Pour le temps, utiliser l'identité de Wald\index{identité de Wald}. Le gain moyen par pari est $0.2$, et le gain final moyen vaut $10P - 5 \approx 3.84$, donc le nombre de paris vaut $3.84/0.2 \approx 19.2$.

\iqwrong{Supposer qu'un taux de réussite de 60~\% signifie une chance de 60~\% de doubler sa mise. L'edge s'accumule sur de nombreux paris, si bien que la probabilité est bien plus élevée.}
\iqwhy{Un petit edge tend vers la certitude dès que la cible est éloignée de nombreux paris. C'est le dimensionnement des positions, et non la seule prévision, qui fixe le risque.}
\end{solution}

\enlargethispage{-3\baselineskip}%
\begin{solution}{iq:iv1:random-walks-and-martingales:3}
On gagne 1~\$ dans 90~\% des séances. Le profit moyen vaut 0, et une séance compte 9 paris en moyenne.

Un jeu équitable reste équitable lorsqu'on l'arrête\index{arrêt optionnel}, tant que l'arrêt est borné comme ici. Donc $0 = P \cdot 1 + (1 - P)(-9)$, ce qui donne $P = 0.9$. Avec $a = 1$ et $b = 9$, le résultat général pour le temps est $ab = 9$.

\iqwrong{Croire qu'un taux de réussite (hit rate) de 90~\% prouve une bonne stratégie. La perte rare est exactement assez grande pour annuler tous les gains. Un taux de réussite élevé ne dit rien ni sur l'équité du jeu, ni sur sa valeur.}
\iqwhy{Vendre de l'assurance ressemble à cela~: beaucoup de petits gains, une grosse perte. Un historique avec 90~\% de taux de réussite peut n'être rien d'autre que la marche entre deux arrêts.}
\end{solution}

\begin{solution}{iq:iv1:random-walks-and-martingales:4}
La probabilité de revenir en 0 après 10 pas vaut $252/1024 \approx 0.246$. Après 100 pas, la distance moyenne vaut environ $7.96$, et $\sqrt{2N/\pi}$, avec $N = 100$ pas, donne $7.98$.

Pour être en 0 après $2n$ pas, exactement $n$ d'entre eux doivent monter, donc $u_{2n} = \binom{2n}{n}/4^n$.

Pour la distance, il existe une formule exacte\index{déplacement absolu moyen}~: $E|S_{2n}| = 2n\,u_{2n}$, ici $100 \cdot 0.0796$. L'approximation vient d'une loi normale de variance $N$, dont la valeur absolue moyenne vaut $\sigma\sqrt{2/\pi}$.

\iqwrong{Répondre 10, l'écart-type. La distance moyenne est plus petite, d'un facteur $\sqrt{2/\pi} \approx 0.8$.}
\end{solution}

\enlargethispage{-3\baselineskip}%
\begin{solution}{iq:iv1:random-walks-and-martingales:5}
Il faut 50 pas en moyenne.

Utiliser l'identité de Wald\index{identité de Wald}~: si une marche à pas indépendants d'espérance $\mu$ s'arrête en un temps d'arrêt $T$ d'espérance finie, alors $E[S_T] = \mu E[T]$. Ici chaque pas a pour espérance $0.2$, et la marche s'arrête exactement en $+10$, car elle avance d'une place à la fois. Donc $E[T] = 10/0.2 = 50$.

Une seconde méthode vérifie l'identité. Soit $x$ le temps moyen pour monter d'un niveau. Le premier pas monte avec probabilité $0.6$. Sinon, il faut regagner deux niveaux~: $x = 1 + 0.4 \cdot 2x$, donc $x = 5$, et dix niveaux prennent 50 pas.

\iqwrong{Utiliser $10/0.6$, en ne comptant que les pas vers le haut. Les pas vers le bas doivent être rattrapés, et c'est la dérive nette qui fixe le rythme.}
\iqfollow{Et lorsque la dérive tend vers zéro~?}{Le temps croît comme $10/\mu$ sans limite~: à dérive nulle il est infini, alors que la marche atteint encore $+10$ avec certitude.}
\end{solution}

\begin{solution}{iq:iv1:random-walks-and-martingales:6}
La marche atteint $+10$ avec une probabilité $(2/3)^{10} \approx 0.0173$, et son point le plus haut culmine en moyenne à 2.

Soit $x$ la probabilité de monter un jour d'un niveau. Par analyse au premier pas\index{analyse au premier pas}, on monte tout de suite, ou on descend et il faut remonter deux niveaux~: $x = 0.4 + 0.6 x^2$. Les racines sont $2/3$ et 1. La marche dérive vers le bas, donc elle n'est pas certaine de remonter, et la petite racine est la bonne. Pour atteindre $+10$, il faut monter d'un niveau dix fois de suite, et chaque niveau est un problème nouveau et indépendant, donc la probabilité vaut $x^{10}$.

Donc le point le plus haut $M$ vérifie $P(M \ge k) = (2/3)^k$, une loi géométrique\index{maximum d'une marche à dérive négative}, et
\[
E[M] = \sum_{k \ge 1} (2/3)^k = \frac{2/3}{1 - 2/3} = 2.
\]
\iqwrong{Prendre la racine $x = 1$ parce que l'équation la permet. Avec une dérive négative, la marche part vers moins l'infini, donc le maximum doit admettre une limite.}
\iqwhy{Cela donne la loi de la perte la plus profonde sous le capital de départ, pour une stratégie à edge positif sur un pari fixe~: vue de l'autre côté, c'est le point le plus haut d'une marche négative. Ce n'est pas le drawdown maximal, qui se mesure depuis chaque nouveau sommet et qui croît d'autant plus qu'on trade longtemps.}
\end{solution}

\begin{solution}{iq:iv1:random-walks-and-martingales:7}
On termine 1~\$ en avance avec une probabilité $1 - 2^{-10} \approx 0.999$. L'espérance du profit vaut 0, et une séance compte environ $1.998$ paris en moyenne.

Les mises $1, 2, 4, \dots, 512$ totalisent 1\,023~\$, donc on peut faire 10 paris. On les perd tous avec une probabilité $2^{-10}$, et alors on perd les 1\,023~\$. Sinon on gagne 1~\$, car tout gain couvre toutes les pertes précédentes plus un. Donc $E[\text{profit}] = (1 - 2^{-10}) - 2^{-10} \cdot 1023 = 0$.

Le nombre de paris est plafonné à 10, et en additionnant $P(\text{le pari } k \text{ a lieu}) = 2^{-(k - 1)}$, on obtient $2 - 2^{-9}$.

\iqwrong{Répondre «~Presque toujours un profit, donc c'est forcément bon~». Par l'arrêt optionnel\index{arrêt optionnel}, aucune règle d'arrêt ne rend un jeu équitable profitable lorsque le capital est fini. La faible probabilité de perte entraîne une perte égale à mille fois le gain.}
\iqfollow{Et avec un capital illimité~?}{On gagne toujours 1~\$, en environ 2 paris, mais le capital nécessaire a une espérance infinie. La méthode exige ce que personne ne possède.}
\end{solution}

\begin{solution}{iq:iv1:random-walks-and-martingales:8}
Cela ne contredit pas l'affirmation, car le théorème exige des hypothèses qui ne sont pas vérifiées ici~: le temps d'arrêt a une espérance infinie, et la marche n'est pas bornée inférieurement. Le temps d'arrêt médian vaut 1, alors que son espérance est infinie.

L'arrêt optionnel\index{arrêt optionnel} reste valable pour un temps borné, ou pour une martingale bornée. Ici aucune des deux hypothèses n'est vérifiée. À tout horizon fixe le jeu est équitable~: sur les trajectoires qui ne se sont pas arrêtées, la position moyenne est bien en dessous de 0, et cela compense les trajectoires arrêtées en $+1$.

Le temps $T$ vaut 1 avec une probabilité $1/2$, ce qui fixe la médiane. Par le principe de réflexion\index{principe de réflexion}, $P(T > 2n) = u_{2n} = \binom{2n}{n}/4^n \approx 1/\sqrt{\pi n}$. Par exemple, $P(T \le 100) \approx 0.92$. La queue de la loi décroît trop lentement pour que la somme $\sum_n P(T > n)$ converge, donc $E[T] = \infty$. Il faudrait survivre à des pertes arbitrairement grandes entre-temps.

\iqwrong{Conclure que la marche ne peut pas être certaine d'atteindre $+1$, parce que le jeu est équitable. Elle l'atteint avec certitude, et l'équité à chaque instant fixe n'empêche pas cela.}
\iqwhy{Une stratégie qui garde les positions perdantes jusqu'à ce qu'elles se redressent est cette marche. Elle affiche un profit à chaque arrêt, et le coût est caché dans le temps illimité et la perte illimitée qu'il faut supporter.}
\end{solution}

\begin{solution}{iq:iv1:random-walks-and-martingales:9}
De 0 à 10, il faut 100 pas en moyenne, et de 4 à 10, il en faut 84.

Utiliser $S_t^2 - t$\index{martingale quadratique}. À l'intérieur de l'intervalle, un pas ajoute en moyenne 1 à $S^2$, donc le processus reste équitable. En 0, le pas forcé fait passer $S^2$ de 0 à 1, ce qui ajoute aussi exactement 1, si bien que le même processus reste équitable au point réfléchissant. L'arrêt en 10 donne $E[S_T^2] - E[T] = k^2$, et donc $E[T] = 100 - k^2$.

Pour $k = 0$ cela donne 100, et pour $k = 4$ cela donne 84.

\iqwrong{Utiliser 24 pas depuis 4, c'est-à-dire le $k(n - k)$ d'une marche avec une sortie en 0. Cette marche-là s'arrêterait en 0, alors que celle-ci rebondit et poursuit sa route.}
\iqfollow{Et en général~?}{Atteindre un niveau $n$ depuis une paroi réfléchissante prend $n^2$ pas. Le temps croît comme le carré de la distance, signature d'une diffusion.}
\end{solution}

\enlargethispage{-4\baselineskip}%
\begin{solution}{iq:iv1:random-walks-and-martingales:10}
La probabilité vaut environ $0.928$, et il faut environ $21.4$ pas en moyenne.

Les pas où la marche reste sur place ne changent rien à la suite de la marche, on peut donc les ignorer. Les pas qui bougent montent avec une probabilité $0.5/0.8 = 5/8$. Donc $r = 3/5$, et la formule de sortie\index{martingale exponentielle} donne
\[
P = \frac{1}{1 + r^5} = \frac{1}{1 + 0.6^5} \approx 0.9279.
\]
En revanche, le temps change. Par l'identité de Wald\index{identité de Wald}, la position finale vaut en moyenne $5P - 5(1 - P) \approx 4.279$. Le pas moyen, zéro compris, vaut $0.5 - 0.3 = 0.2$. Donc le temps vaut $4.279/0.2 \approx 21.4$.

\iqwrong{Exclure les pas sur place de la durée, ou les laisser modifier la probabilité de sortie. Ils sont sans effet sur le point de sortie, mais pas sur le temps mis pour y parvenir.}
\iqwhy{Les données de mise à jour des cotations comportent souvent beaucoup de pas inchangés. Les éliminer modifie l'horloge sans changer la direction, et il est plus propre de séparer ces deux effets.}
\end{solution}

\begin{solution}{iq:iv1:random-walks-and-martingales:11}
La cible est atteinte en premier avec une probabilité d'environ $0.868$, la transaction dure environ 160 ticks et rapporte en moyenne environ $16.03$ ticks. Une marche équitable donnerait une chance de $1/3$ et un profit espéré nul.

Avec $r = 0.45/0.55 = 9/11$, la formule de sortie\index{martingale exponentielle} partant de 10 ticks au-dessus du stop, l'écart entre le stop et l'objectif étant de 30 ticks, vaut
\[
P = \frac{1 - r^{10}}{1 - r^{30}} \approx 0.8677.
\]
Le profit espéré vaut $20P - 10(1 - P) = 30P - 10 \approx 16.03$ ticks. Chaque tick rapporte $0.1$ en moyenne, donc l'identité de Wald\index{identité de Wald} donne $16.03/0.1 \approx 160$ ticks.

Avec une marche équitable, $P = 10/30 = 1/3$, et le profit est nul quels que soient le stop et l'objectif retenus. Les stops et les objectifs modifient la forme du résultat, pas sa moyenne. L'edge d'un trader provient des probabilités de tick, et ces sorties ne font que le récolter.

\iqwrong{«~Un objectif double du stop, donc il suffit de gagner un tiers des transactions.~» C'est une affirmation sur le seuil de rentabilité, pas sur un edge. Elle peut être vraie d'une marche équitable, et elle ne rapporte rien.}
\iqfollow{Combien de temps met une marche équitable~?}{À partir de la martingale $S_t^2 - t$, le temps d'une marche équitable est le produit des distances aux deux sorties~: $10 \cdot 20 = 200$ ticks en moyenne. C'est plus long qu'avec la dérive, car la marche ne reçoit aucune aide pour se diriger vers l'objectif.}
\end{solution}

\begin{solution}{iq:iv1:random-walks-and-martingales:12}
Environ $1.9 \times 10^{-9}$ pour des mises de 1~\$, $0.134$ pour 10~\$, et $0.669$ pour 50~\$.

On change d'échelle la marche\index{risque de ruine}. Avec des mises de $s$, le capital vaut $100/s$ mises, et la probabilité de les perdre un jour vaut $r^{100/s}$, avec $r = q/p = 0.45/0.55 = 9/11$~: c'est la probabilité qu'une marche de dérive positive descende un jour d'autant de niveaux. Il faut donc $r^{100}$, $r^{10}$ et $r^2$.

On note ce qui ne change pas~: chaque mise a le même edge en pourcentage. L'espérance de gain par dollar misé est la même, mais la ruine passe d'un risque quasi nul à un risque probable à mesure que la mise croît, car le capital mesuré en mises est ce qui protège.

\iqwrong{Penser que la ruine ne dépend que de l'edge. Avec un edge positif, on fait tout de même faillite si la mise est trop grande, et une pièce équilibrée fait faillite à coup sûr, quelle que soit la mise.}
\iqwhy{Le risque de ruine décroît exponentiellement avec le capital mesuré en mises. D'où la taille de la mise compte plus que presque tout le reste pour la survie~: doubler la mise remplace la probabilité de ruine par sa racine carrée, ce qui est bien plus grand.}
\end{solution}

\begin{solution}{iq:iv1:random-walks-and-martingales:13}
Elle reste strictement au-dessus de 0 pendant 20 pas avec probabilité $\tfrac12 u_{20} \approx 0.0881$. Elle revient à 0 pour la première fois aux pas 2, 4 et 6 avec probabilités $1/2$, $1/8$ et $1/16$.

Ici $u_{2n} = \binom{2n}{n}/4^n$ est la probabilité d'être en 0 à l'instant $2n$. Par le théorème du scrutin\index{théorème du scrutin}\index{principe de réflexion}, les chemins qui restent strictement positifs pendant $2n$ pas sont exactement la moitié de ceux qui finissent en 0, ce qui donne $\tfrac12 u_{2n}$. Donc $u_{20} \approx 0.176$ est la probabilité de ne jamais revenir à 0 en 20 pas, dans un sens comme dans l'autre.

Pour le premier retour, les mêmes idées donnent $f_{2n} = u_{2n}/(2n - 1)$\index{temps du premier retour}. Donc $f_2 = 1/2$, $f_4 = (3/8)/3 = 1/8$ et $f_6 = (5/16)/5 = 1/16$. L'énumération de tous les chemins le confirme. La marche revient à 0 avec certitude, mais l'espérance du temps de retour est infinie.

\iqwrong{Utiliser $2^{-20}$, ce qui traite «~rester au-dessus de 0~» comme si chaque pas devait monter. La plupart des chemins qui restent au-dessus de 0 montent et descendent beaucoup.}
\iqwhy{Un processus équitable peut rester d'un côté pendant longtemps avec une probabilité assez élevée. Une série de bons P\&L qui dure 20 périodes n'a qu'environ une chance sur onze.}
\end{solution}

\begin{solution}{iq:iv1:random-walks-and-martingales:14}
La probabilité vaut $(1/3)/(8/9) = 3/8$. La réponse repose seulement sur le fait qu'une croyance correcte est une martingale\index{croyance comme martingale}.

Si la croyance est la bonne probabilité conditionnelle, son espérance pour la valeur suivante égale sa valeur actuelle~: sinon on croirait déjà autre chose. On arrête donc au premier instant où elle atteint $8/9$. Les cotes évoluent par facteurs de 2, donc elle tombe exactement sur $8/9$. Si elle n'y parvient jamais, elle tend finalement vers 0, car elle converge vers la vérité. Le théorème d'arrêt optionnel\index{arrêt optionnel} donne
\[
\tfrac13 = P \cdot \tfrac89 + (1 - P) \cdot 0, \quad\text{donc}\quad P = \tfrac38.
\]
Ni la fiabilité de l'information ni le nombre de jours n'apparaissent, car une croyance correcte évolue de sorte que ce résultat tient dans tous les cas. Des cotes allant de $1 : 2$ à $8 : 1$ en font une marche propre, et une simulation avec un état vrai tiré dans la loi a priori donne le même résultat.

\iqwrong{Répondre $8/9$, ou penser que, parce que l'événement est peu probable, la croyance ne peut pas monter. Une croyance à $1/3$ montera souvent puis retombera, car un événement à $1/3$ échoue deux fois sur trois.}
\iqwhy{Le prix d'un contrat binaire est une croyance de ce type. À 30 cents, il peut atteindre 90 au plus un tiers du temps, quoi qu'il arrive.}
\end{solution}

\begin{solution}{iq:iv1:random-walks-and-martingales:15}
En moyenne, il leur faut 16 pas.

Ne suivre que leur écart\index{agrégation d'états}. L'écart part de 4 sur un anneau de 12. À chaque pas, chaque jeton fait $\pm 1$, donc l'écart varie de $-2$ avec probabilité $1/4$, de $+2$ avec probabilité $1/4$, et de 0 avec probabilité $1/2$. Comme les écarts restent pairs, on mesure en unités de 2~: l'écart vit sur un anneau de 6 positions, part à 2 du point de rencontre, et c'est une marche qui reste sur place la moitié du temps.

Une marche sur un anneau partant à distance 2 sur 6 est la même chose que la ruine du joueur\index{ruine du joueur} sur une droite partant de 2 entre 0 et 6, elle prend donc $2 \cdot 4 = 8$ pas. La moitié du temps rien ne bouge, ce qui le porte à 16 pas.

\iqwrong{Utiliser $4 \cdot 8 = 32$, en traitant l'écart comme une marche simple sur 12 positions. Les deux jetons bougent ensemble, ce qui fait avancer l'écart de deux places à la fois, et ne le fait pas bouger du tout la moitié du temps.}
\iqfollow{Et s'ils partent à 3 d'écart~?}{Ils ne se rencontrent jamais, car la parité de l'écart ne change jamais quand les deux bougent.}
\end{solution}

\begin{solution}{iq:iv1:random-walks-and-martingales:16}
La dernière visite en 0 a le plus de chances de se produire tout au début ou tout à la fin~: à l'instant 0 et à l'instant 20, chacun avec une probabilité d'environ $0.176$, et l'instant 10 est le moins probable, d'environ $0.061$. À la limite, la dernière visite tombe dans le premier quart avec probabilité $1/3$.

La probabilité que la dernière visite ait lieu à l'instant $2k$ est la probabilité d'être en 0 à cet instant, multipliée par la probabilité de ne pas revenir pendant le temps restant~; cette seconde probabilité est aussi une valeur de $u$~: $P = u_{2k}\, u_{2n - 2k}$. La loi est donc symétrique, et elle prend de grandes valeurs aux deux extrémités, car $u$ décroît avec le temps. Pour $n = 10$~: $u_{20} \approx 0.176$ aux deux extrémités, et $u_{10}^2 \approx 0.061$ au milieu.

À la limite, cela devient la loi de l'arcsinus\index{loi de l'arcsinus}~: la fraction $x$ de la durée au bout de laquelle a lieu la dernière visite vérifie
\[
P(\text{dernière visite avant } x) = \frac{2}{\pi} \arcsin\sqrt{x},
\]
ce qui donne $1/3$ pour $x = 1/4$. Donc avec probabilité $2/3$, la dernière visite tombe dans le dernier quart ou dans le premier quart.

\iqwrong{Penser que la dernière visite se situe près du milieu de la durée, ou qu'une marche équitable doit passer environ la moitié du temps de chaque côté. En vérité, une marche équitable reste en général d'un seul côté pendant la plus grande partie de sa vie.}
\iqwhy{Une stratégie plate avec une longue série gagnante ou une longue série perdante est exactement ce que produit le hasard. La loi de l'arcsinus met en garde contre l'inférence de changements de régime à partir de longues séries de même signe dans les P\&L cumulés.}
\end{solution}

\begin{solution}{iq:iv1:random-walks-and-martingales:17}
Le premier joueur n'a plus de pièces après 300 tours en moyenne, et un joueur détient les 60 pièces après 1\,100 tours en moyenne. Les formules générales sont $3abc/(a + b + c)$ et $ab + bc + ca$.

Chercher des quantités qui décroissent régulièrement\index{martingale issue d'un produit}. Quand le joueur $x$ donne une pièce au joueur $y$, $z$ étant le troisième joueur, le produit $abc$ varie de $z(x - y - 1)$. Sur les six couples ordonnés, les parties en $x - y$ se compensent, et la variation moyenne vaut $-(a + b + c)/3$, une constante, car le total reste fixe. Donc $abc + t(a + b + c)/3$ est une martingale. Elle s'arrête au moment où un joueur n'a plus de pièces, quand le produit vaut 0, donc
\[
abc = E[T] \cdot \frac{a + b + c}{3}, \quad\text{donc}\quad E[T] = \frac{3abc}{a + b + c} = \frac{3 \cdot 6000}{60} = 300.
\]
Pour la seconde partie, utiliser $ab + bc + ca$. Tout transfert la modifie de $x - y - 1$, et sa moyenne vaut $-1$, pour trois joueurs comme pour deux. Elle baisse donc en moyenne d'une unité par tour, et elle atteint 0 lorsqu'un joueur détient tout. Donc $E[T] = 200 + 600 + 300 = 1100$.

Résoudre le système linéaire complet confirme les deux formules pour de petits départs, comme 3 et 11 à partir de $(1, 2, 3)$.

\iqwrong{Vouloir écrire la récurrence sur l'état. Il y a beaucoup d'états, et l'intérêt de la question est qu'une bonne quantité donne la réponse en une ligne.}
\iqfollow{Et avec deux joueurs, de fortunes $a$ et $b$~?}{Le produit $ab$ baisse d'une unité par tour, ce qui retrouve le $ab$ de la ruine du joueur.}
\end{solution}

\begin{solution}{iq:iv1:random-walks-and-martingales:18}
La borne supérieure vaut $2/3$, la borne inférieure obtenue à long terme vaut $1/4$, et un calcul direct donne environ $0.44$.

\begin{sloppypar}Martingale~: si le rouge représente la fraction $f = r/(r + b)$, la fraction suivante a pour espérance
\[
f \cdot \frac{r + 1}{r + b + 1} + (1 - f) \cdot \frac{r}{r + b + 1} = \frac{r^2 + r + br}{(r + b)(r + b + 1)} = f.
\]
L'inégalité maximale de Doob\index{inégalité maximale de Doob} pour une martingale positive donne $P(\sup \ge c) \le f_0/c$, ici $(1/2)/(3/4) = 2/3$. Il ne s'agit que d'un arrêt optionnel, la fraction valant au moins $c$ sur cet événement.\end{sloppypar}

Borne inférieure~: la fraction converge\index{urne de Pólya}, et sa limite est uniforme sur $[0, 1]$. Elle finit donc au-dessus de $3/4$ avec probabilité $1/4$, et cela est compris dans les instants où elle l'a atteint. Une minoration plus fine et simple vient de deux tirages rouges consécutifs, qui donnent déjà $3/4$~: $\tfrac12 \cdot \tfrac23 = 1/3$.

Une programmation dynamique sur les effectifs de l'urne jusqu'à 4\,000 boules donne environ $0.44$, ce que confirme une simulation. La fraction dépasse souvent $3/4$ tôt, lorsque peu de boules sont dans l'urne, puis redescend. C'est pourquoi la probabilité pour la limite est bien inférieure à celle de l'atteindre à un moment.

\iqwrong{Prendre la borne pour la réponse. La borne de Doob n'est atteinte que lorsque la trajectoire touche exactement le niveau et finit en 0 sinon, et ici la fraction ne peut pas tomber à 0.}
\end{solution}

\begin{solution}{iq:iv1:random-walks-and-martingales:19}
La moyenne vaut environ $7.50$. La valeur continue $\sqrt{2n/\pi}$ vaut $7.98$, et la correction de $\tfrac12$ donne $7.48$.

Le principe de réflexion\index{principe de réflexion} donne, pour $a \ge 1$,
\[
P(M \ge a) = P(S_n \ge a) + P(S_n > a),
\]
car chaque trajectoire qui atteint $a$ et finit en dessous peut être réfléchie, après son premier contact avec $a$, en une trajectoire finissant au-dessus. Alors $E[M] = \sum_{a \ge 1} P(M \ge a)$, et l'ajout des queues binomiales donne $7.50$, ce qu'une énumération complète confirme sur des marches plus courtes. Les trajectoires qui touchent exactement $a$ ne sont comptées qu'une fois, et non deux, et c'est là que naît l'écart avec la formule brownienne.

Dans la limite continue\index{maximum du mouvement brownien}, $M$ a la même loi que $|W_n|$, dont la moyenne vaut $\sqrt{2n/\pi}$. La marche discrète progresse par pas d'une unité entière, elle reste donc environ un demi-pas en deçà de la limite, et la correction de $\tfrac12$ en provient.

\iqwrong{Répondre 10, ou doubler $E[S_n^+]$ sans correction. L'écart type de la marche n'est pas la moyenne de son maximum. Sur des marches courtes et pour des barrières, la correction du demi-pas est non négligeable.}
\iqwhy{Une barrière à surveillance discrète se trouve en pratique plus loin qu'une barrière continue de même niveau. La correction usuelle décale la barrière d'une fraction de l'écart type d'un pas, ce qui traduit ce même fait.}
\end{solution}

\begin{solution}{iq:iv1:random-walks-and-martingales:20}
La variance vaut 1\,200, donc l'écart type est d'environ $34.6$ pas.

Chaque pas a une moyenne $\mu = 0.2$ et une variance $\sigma^2 = 1 - \mu^2 = 0.96$. Le processus $(S_t - \mu t)^2 - \sigma^2 t$ est une martingale\index{identité de Wald du second ordre}. On l'arrête en $T$, lorsque $S_T = 10$ exactement~:
\[
E[(10 - \mu T)^2] = \sigma^2 E[T].
\]
Comme $E[T] = 10/\mu$, le membre de gauche vaut $\mu^2\,\mathrm{Var}(T)$. Donc
\[
\mathrm{Var}(T) = \frac{\sigma^2 E[T]}{\mu^2} = \frac{10\,\sigma^2}{\mu^3} = \frac{10 \cdot 0.96}{0.008} = 1200.
\]
Une simulation de 20\,000 marches donne le même résultat. Comme on s'y attend, la variance croît linéairement avec la distance, car chaque niveau correspond à une attente distincte et indépendante.

\iqwrong{Utiliser une variance de $0.96 \cdot 50 = 48$, la variance de la position à un instant fixé. L'attente est bien plus variable, puisque la variance est divisée par $\mu^3$, et une petite dérive donne une dispersion énorme.}
\iqwhy{Pour une stratégie à faible avantage, le temps pour atteindre une cible a une loi très dispersée~: ici un écart type de 70~\% de la moyenne. On doit s'attendre à de longues périodes sans progrès, même lorsque l'avantage est réel.}
\end{solution}

\begin{solution}{iq:iv1:random-walks-and-martingales:21}
La marche revient avec certitude en une et deux dimensions, et avec une probabilité d'environ $0.34$ en trois dimensions. C'est le théorème de Pólya\index{théorème de Pólya}.

Le nombre de visites au point de départ est un compte géométrique~: avec une probabilité de retour $R$, sa moyenne vaut $U = 1/(1 - R)$. La marche revient donc avec certitude exactement lorsque $U = \sum_n u_n$ diverge.

En dimension un, $u_{2n} = \binom{2n}{n}/4^n \sim 1/\sqrt{\pi n}$, et la somme diverge. En dimension deux, la rotation du réseau de 45 degrés transforme la marche en deux marches unidimensionnelles indépendantes, donc $u_{2n}$ est le carré de la valeur unidimensionnelle, soit environ $1/(\pi n)$, et la somme diverge encore, tout juste. En dimension trois, $u_{2n}$ décroît comme environ $n^{-3/2}$, et la somme converge.

La valeur utilise la forme de Fourier de $U$, la moyenne de $1/(1 - \tfrac13(\cos a + \cos b + \cos c))$ sur le cube des angles, qui vaut environ $1.516$. Donc $R = 1 - 1/1.516 \approx 0.34$. Une simulation directe donne le même résultat.

\iqwrong{Penser qu'une marche doit forcément revenir, puisqu'elle dispose d'un temps infini. En dimension assez élevée, il y a trop d'endroits où aller.}
\iqfollow{Quelles sont les cotes en dimension quatre~?}{Plus faibles encore, environ $0.19$ par le même raisonnement. La probabilité de retour décroît régulièrement avec la dimension.}
\end{solution}

\begin{solution}{iq:iv1:random-walks-and-martingales:22}
On s'arrête en haut avec une probabilité d'environ $0.326$, en environ $25.9$ pas. Un tiers est proche, mais la sortie supérieure peut dépasser le niveau et atteindre 11.

La marche a un pas moyen nul, et elle descend d'une place à la fois, donc elle atterrit exactement sur $-5$. En montant, elle atterrit en revanche sur 10 ou 11. Soient $a$ et $b$ les probabilités de s'arrêter en 10 et en 11. L'arrêt optionnel\index{arrêt optionnel} donne $10a + 11b - 5(1 - a - b) = 0$, soit $15a + 16b = 5$. Cela fait deux inconnues, il faut donc un fait supplémentaire, et un système linéaire sur les positions le fournit, avec $a \approx 0.218$ et $b \approx 0.108$. Leur total vaut $0.326$.

Pour le temps, chaque pas a une variance 2 (c'est $\tfrac13 \cdot 4 + \tfrac23 \cdot 1$), donc $S_t^2 - 2t$ est une martingale. Alors $E[T] = \tfrac12 E[S_T^2] = \tfrac12(100a + 121b + 25(1 - a - b)) \approx 25.9$.

Si la marche atterrissait exactement sur 10, comme le fait une marche à pas unitaires, on obtiendrait $5/15 = 1/3$~; avec la variance de pas de 2 de cette marche, le temps vaut $5 \cdot 10/2 = 25$. (Avec des pas unitaires, de variance 1, il faut 50 pas.) Le dépassement signifie que le sommet se trouve en moyenne un peu plus loin que 10. On sort donc un peu moins souvent par là.

\iqwrong{Ignorer le dépassement parce que la réponse semble proche. Une réponse d'un tiers ne s'écarte que d'environ 2~\%, mais elle est fausse, et avec des sauts plus grands elle peut l'être beaucoup plus. L'identité de Wald exige la position à l'arrêt, et ce n'est pas la barrière lorsque les sauts peuvent la franchir.}
\iqwhy{Un ordre stop exécuté à travers un gap est un dépassement. Le niveau réellement atteint n'est pas la barrière, et un saut du cours est ce qui met en échec le calcul classique.}
\end{solution}

\enlargethispage{-1\baselineskip}%
\begin{solution}{iq:iv1:random-walks-and-martingales:23}
La variance vaut $\tfrac23 n^2(n^2 - 1)$, soit 6\,600 pour $n = 10$, donc l'écart type est d'environ 81 pour une moyenne de 100.

On utilise une martingale de puissance quatrième\index{martingale de puissance quatrième}. Sur un pas, $E[(S + 1)^4 + (S - 1)^4]/2 = S^4 + 6S^2 + 1$. En tenant compte de ce terme, on vérifie que
\[
M_t = S_t^4 - 6tS_t^2 + 3t^2 + 2t
\]
reste équitable sur un pas. On l'arrête en $T$, avec $|S_T| = n$~:
\[
0 = n^4 - 6n^2 E[T] + 3E[T^2] + 2E[T].
\]
Comme $E[T] = n^2$, cela donne $3E[T^2] = 5n^4 - 2n^2$. Donc $E[T^2] = \tfrac13(5n^4 - 2n^2)$, et en retranchant $n^4$,
\[
\mathrm{Var}(T) = \frac{2n^4 - 2n^2}{3} = \frac23 n^2(n^2 - 1).
\]
Vérification~: pour $n = 1$, on sort après un pas avec certitude, et la formule donne 0. Les équations de l'analyse au premier pas sur le second moment donnent la même valeur pour $n$ jusqu'à 10, et une simulation la confirme.

\iqwrong{Supposer que l'écart relatif diminue quand $n$ croît. Le rapport entre écart type et moyenne tend vers $\sqrt{2/3}$, donc l'attente reste très variable.}
\iqfollow{Pourquoi ce rapport est-il fixe~?}{Par un argument de changement d'échelle. Une marche à l'échelle $n$ se comporte comme un mouvement brownien sur $[-1, 1]$, le temps étant rééchelonné par $n^2$, donc tous les moments de $T/n^2$ ont une limite.}
\end{solution}

\begin{solution}{iq:iv1:random-walks-and-martingales:24}
En moyenne, elle visite $+k$ exactement une fois, pour chaque $k$, y compris 100.

Pour la preuve, réfléchissons à la manière dont la marche atteint $k$. Elle fait d'abord un pas vers $+1$, avec probabilité $1/2$. Elle atteint ensuite $k$ avant 0, par la ruine du joueur\index{ruine du joueur}, avec probabilité $1/k$. Comptons maintenant les visites en $k$ avant 0. Depuis $k$, un pas va dans un sens ou dans l'autre. En montant, elle revient toujours. En descendant, elle atteint 0 avant de revenir avec probabilité $1/k$. Chaque visite a donc une probabilité d'échappement de $\tfrac12 \cdot \tfrac1k$, et le nombre de visites est géométrique, de moyenne $2k$. Au total~:
\[
\tfrac12 \cdot \tfrac1k \cdot 2k = 1.
\]
C'est ce à quoi l'on s'attend d'après la mesure stationnaire\index{mesure stationnaire}. La mesure de comptage est stationnaire pour la marche, et entre deux retours sur un état quelconque, le nombre moyen de visites dans un autre état est le rapport des masses stationnaires, ici $1$. Sur un anneau fini, cela vaut exactement, pour tout $k$, comme le confirme la résolution du système linéaire.

Ce n'est pas absurde, puisque la marche atteint $+100$ seulement environ une fois sur 200 excursions. Mais quand elle y parvient, elle y effectue environ 200 visites. Les excursions rares compensent cela, et comme chaque niveau apporte $1$, le temps total est infini.

\iqwrong{Affirmer que le niveau lointain doit être visité bien moins souvent, car la plupart des retours sont rapides. La plupart des excursions n'atteignent jamais ce niveau, mais les rares qui y parviennent y restent longtemps.}
\iqwhy{Les excursions rares et longues dominent les moyennes. Une moyenne sur de nombreuses excursions courtes les manque dans un petit échantillon, et elle peut même être indéfinie. C'est une mise en garde sur ce qu'un petit backtest apprend d'un processus à queue épaisse.}
\end{solution}

\begin{solution}{iq:iv1:random-walks-and-martingales:25}
Avec une mise de 1~\$, la probabilité est d'environ $6.2 \times 10^{-14}$, et le jeu audacieux donne $0.16$, qui est le meilleur possible. Pour une probabilité de gain de $0.6$, les petites mises sont les meilleures, avec une probabilité proche de 1. Pour $0.5$, chaque plan donne $1/4$.

Miser 1~\$ revient à une ruine du joueur\index{ruine du joueur} avec $r = q/p = 3/2$, en partant de 25 avec un objectif de 100~:
\[
P = \frac{1.5^{25} - 1}{1.5^{100} - 1} \approx 6.2 \times 10^{-14}.
\]
Le jeu audacieux\index{jeu audacieux} mise la totalité des 25~\$, puis 50~\$~: deux gains d'affilée, de probabilité $0.4^2 = 0.16$. L'itération sur la valeur, menée sur toutes les mises en dollars entiers, donne aussi $0.16$ comme meilleure valeur.

Pourquoi~? Chaque pari à $p < 1/2$ coûte, en espérance, proportionnellement au montant misé. Les petites mises imposent de miser un total important avant d'atteindre l'objectif, et chaque dollar misé perd de l'argent en moyenne. Le jeu audacieux engage le plus petit total misé permettant d'atteindre l'objectif.

À $p = 0.6$, la logique s'inverse~: chaque pari rapporte en espérance, il faut donc miser beaucoup au total. La probabilité avec de petites mises est $(1 - r^{25})/(1 - r^{100})$ avec $r = 2/3$, soit environ $0.99996$, contre $0.36$ pour le jeu audacieux. Pour $p = 0.5$, la richesse est une martingale, donc l'arrêt optionnel donne une probabilité de $25/100$ pour toute stratégie.

\iqwrong{Affirmer que les petites mises réduisent le risque, donc qu'elles doivent être plus sûres ici. Les petites mises réduisent la variance, et lorsqu'il faut atteindre un objectif face à un edge négatif, c'est la variance que l'on achète.}
\iqwhy{Cela s'applique à un trader en retard sur un objectif, sans véritable edge~: le seul moyen de l'atteindre est d'augmenter la taille des positions, et c'est précisément pour cela que les limites de risque existent. Inversement, un edge positif est détruit si l'on mise comme s'il s'agissait d'une loterie.}
\end{solution}
